\section{Results}
\label{sec:results}

% TODO(results/overview): none needed beyond the subsection order; every number below comes from the 05_report.py output, never typed by hand.

\subsection{What the Published Inputs Were}
\label{sec:results:inputs}

Table~\ref{tab:lineage} compares the inputs recovered from the Phase 4A batch file with each candidate series, and Figures~\ref{fig:window} and~\ref{fig:series} show one window and the 2008--2026 history. \tbd{}
% TODO(results/inputs): (i) 1,476 windows recovered with ends 2020-02-13 to 2025-12-31 and the January calendar (no 2023-07-12, no
% 2024-03-11); (ii) agreement of the daily values with each candidate series by year: equal to new_full at the 0.01B precision of
% the prompt through about 2023-09 and equal to call-only GEX from about 2023-10 through 2025-12; (iii) the old_bs archive is put-less
% for the same dates and has holes, so it cannot stand in for a clean "old convention"; (iv) the inputs of the November 2025 Phase 3
% and 4 runs (252 values each; 2024 gross magnitude with a negative sign imposed, negative on 242 days, ten placeholder rows; 2020
% magnitude only, positive on every day; ISO dates in the prompts). Source: the series table (series jan2026_phase4a against the
% others; p3_2024 and p4_2020 as described in Section 3.2).

\begin{table}[H]
\centering
\caption{Agreement of the recovered published inputs with each candidate series, by year (share of daily values equal at the 0.01B precision of the prompt, in percent).}
\label{tab:lineage}
\footnotesize
\begin{tabular}{@{}lrrrr@{}}
\toprule
Year & Days & \ident{new_full} & \ident{call_only} & \ident{old_bs} \\
\midrule
2020 & \tbd{} & \tbd{} & \tbd{} & \tbd{} \\
2021 & \tbd{} & \tbd{} & \tbd{} & \tbd{} \\
2022 & \tbd{} & \tbd{} & \tbd{} & \tbd{} \\
2023 & \tbd{} & \tbd{} & \tbd{} & \tbd{} \\
2024 & \tbd{} & \tbd{} & \tbd{} & \tbd{} \\
2025 & \tbd{} & \tbd{} & \tbd{} & \tbd{} \\
\bottomrule
\end{tabular}
\end{table}

\begin{figure}[H]
\centering
\resultfigure{\figWindow}
\caption{One SPY 30-day window under four conventions (\ident{jan2026_phase4a}, \ident{call_only}, \ident{new_ex0dte}, \ident{new_full}), window ending \tbd{}. Dashed lines mark the $\pm$\$5 billion magnitude threshold; the table under the bars gives the three criteria and the verdict for each convention. \textbf{Read this figure as:} \tbd{}}
\label{fig:window}
\end{figure}

\begin{figure}[H]
\centering
\resultfigure{\figSeries}
\caption{SPY, 2008--2026. (a) Yearly mean of $|\mathrm{GEX}|$ under each series, with the \$5 billion magnitude criterion for scale. (b) Daily sign agreement between pairs of series, by year. \textbf{Read this figure as:} \tbd{}}
\label{fig:series}
\end{figure}

\subsection{Classifier Replay by Year}
\label{sec:results:replay}

Table~\ref{tab:replay} and Figure~\ref{fig:replay} give the detection rate of the deterministic classifier by year for each series next to the published o4-mini rates, and Figure~\ref{fig:thresholds} shows how far the replay depends on the thresholds. \tbd{}
% TODO(results/replay): data audit, to be verified against the final tables: the new_full replay gives 10.3 % (2024) and 13.6 % (2025)
% against 100 % and 100 % published, and on jan2026_phase4a the replay reproduces the published labels in 1,412 of 1,412 parsed
% windows. Report Wilson intervals and, next to any p-value, the block bootstrap of Section 3.8. The threshold grid has 45
% configurations (persistence 60-80 %, magnitude 3/5/7 billion, flips 3/5/7).

\begin{table}[H]
\centering
\caption{Detection rate of the deterministic classifier by year, SPY, in percent (windows detected out of windows), next to the published o4-mini rate.}
\label{tab:replay}
\footnotesize
\begin{tabular}{@{}lrrrrrrr@{}}
\toprule
Year & Windows & Published & \ident{jan2026_phase4a} & \ident{call_only} & \ident{old_bs} & \ident{new_ex0dte} & \ident{new_full} \\
\midrule
2020 & \tbd{} & \tbd{} & \tbd{} & \tbd{} & \tbd{} & \tbd{} & \tbd{} \\
2021 & \tbd{} & \tbd{} & \tbd{} & \tbd{} & \tbd{} & \tbd{} & \tbd{} \\
2022 & \tbd{} & \tbd{} & \tbd{} & \tbd{} & \tbd{} & \tbd{} & \tbd{} \\
2023 & \tbd{} & \tbd{} & \tbd{} & \tbd{} & \tbd{} & \tbd{} & \tbd{} \\
2024 & \tbd{} & \tbd{} & \tbd{} & \tbd{} & \tbd{} & \tbd{} & \tbd{} \\
2025 & \tbd{} & \tbd{} & \tbd{} & \tbd{} & \tbd{} & \tbd{} & \tbd{} \\
\bottomrule
\end{tabular}
\end{table}

\begin{figure}[H]
\centering
\resultfigure{\figReplay}
\caption{Detection rate by year, SPY, 2020--2025: the published o4-mini labels and the deterministic rule replayed on the published inputs (\ident{jan2026_phase4a}), the call leg only (\ident{call_only}), the pre-rebuild archive (\ident{old_bs}) and the rebuilt table without and with same-day expiries (\ident{new_ex0dte}, \ident{new_full}), with 95\% Wilson intervals. \textbf{Read this figure as:} \tbd{}}
\label{fig:replay}
\end{figure}

\begin{figure}[H]
\centering
\resultfigure{\figThresholds}
\caption{Sensitivity of the 2024--2025 detection rate to the persistence, magnitude and flip thresholds (45 threshold sets: one panel per magnitude and flip limit, persistence on the horizontal axis) for the published inputs (\ident{jan2026_phase4a}) and the rebuilt series (\ident{new_full}); the panel with the heavy frame holds the published thresholds. \textbf{Read this figure as:} \tbd{}}
\label{fig:thresholds}
\end{figure}

\subsection{Model Agreement on the Published Inputs}
\label{sec:results:bridge}

Table~\ref{tab:bridge} and Figure~\ref{fig:bridge} compare gpt-6-luna with the published o4-mini answers on identical inputs. \tbd{}
% TODO(results/bridge): arm bridge_jan against r1_o4mini on (symbol, window_end); gate: Cohen's kappa of at least about 0.6, otherwise
% discuss before reading the other arms. On the exact inputs the o4-mini-reported mean magnitude equals the computed mean (median
% ratio 1.0 in the data audit), which closes the "about 1.25 times" discrepancy seen against the archive; report it here.
% State how the 64 re-parsed answers are handled and give the result with and without them.

\begin{table}[H]
\centering
\caption{Bridge to the published result: gpt-6-luna versus o4-mini on the exact Phase 4A inputs.}
\label{tab:bridge}
\footnotesize
\begin{tabular}{@{}lrrr@{}}
\toprule
Statistic & o4-mini & gpt-6-luna & Both \\
\midrule
Windows with a parsed answer & \tbd{} & \tbd{} & \tbd{} \\
Detection rate (\%) & \tbd{} & \tbd{} & -- \\
Agreement with the classifier (\%) & \tbd{} & \tbd{} & -- \\
Agreement between models (\%) & -- & -- & \tbd{} \\
Cohen's kappa [95\% block-bootstrap interval] & -- & -- & \tbd{} \\
McNemar exact p (non-overlapping subsample) & -- & -- & \tbd{} \\
Block-bootstrap p (calendar quarters) & -- & -- & \tbd{} \\
Median reported / computed mean $|\mathrm{GEX}|$ & \tbd{} & \tbd{} & -- \\
\bottomrule
\end{tabular}
\end{table}

\begin{figure}[H]
\centering
\resultfigure{\figBridge}
\caption{Bridge to the published result on the Phase 4A windows. (a) Confusion matrix of gpt-6-luna against the o4-mini labels, with the overall agreement and Cohen's kappa. (b), (c) Agreement and kappa by year, with the test-retest of gpt-6-luna and the agreement gate as dashed references. \textbf{Read this figure as:} \tbd{}}
\label{fig:bridge}
\end{figure}

\subsection{gpt-6-luna on the Complete Series}
\label{sec:results:corrected}

Figure~\ref{fig:corrected} shows the detection rate of gpt-6-luna by year on each SPY series, from the published inputs to the complete series, next to the classifier. \tbd{}
% TODO(results/corrected): arms full_spy and ex0dte_spy, 2008-2026, by year with Wilson intervals and block-bootstrap intervals; list
% windows where model and classifier disagree; give parse statistics (finish_reason, retries) in one sentence.

\begin{figure}[H]
\centering
\resultfigure{\figCorrected}
\caption{gpt-6-luna detection rate by year, SPY: one panel per arm, from the published inputs (\ident{bridge_jan}) to the rebuilt series with and without same-day expiries (\ident{full_spy}, \ident{ex0dte_spy}) and, if it was run, the pre-rebuild archive (\ident{bridge_old}). Dots are the model with 95\% Wilson intervals, dashes the classifier on the windows the model answered, and the grey band is 2026. \textbf{Read this figure as:} \tbd{}}
\label{fig:corrected}
\end{figure}

\subsection{0DTE Ablation}
\label{sec:results:ablation}

Table~\ref{tab:ablation} separates the contribution of same-day expiries by comparing \ident{new_full} with \ident{new_ex0dte}. \tbd{}
% TODO(results/ablation): share of mean |GEX| carried by same-day expiries by year; paired comparison per window (McNemar, block
% bootstrap). Data audit example for Figure 2: SPY 2024-03-15 is +14.61B on the published inputs and as call-only GEX, -4.50B ex-0DTE and -15.35B
% on the full series (the archive value, +12.62B, is put-less too, so do not read the first step as a pure gamma-model effect).

\begin{table}[H]
\centering
\caption{0DTE ablation, SPY, by year: share of mean $|\mathrm{GEX}|$ carried by same-day expiries (percent), and detection rates (percent) of the classifier and of gpt-6-luna on \ident{new_full} and \ident{new_ex0dte}. The McNemar p-values are exact tests on the non-overlapping subsample (Section~\ref{sec:methods:statistics}); the block-bootstrap values accompany them in \ident{results/stats.yaml}.}
\label{tab:ablation}
\footnotesize
\begin{tabular}{@{}lrrrrrr@{}}
\toprule
 & 0DTE share & \multicolumn{2}{c}{Classifier} & \multicolumn{2}{c}{gpt-6-luna} & McNemar \\
\cmidrule(lr){3-4}\cmidrule(lr){5-6}
Year & of $|\mathrm{GEX}|$ & full & ex-0DTE & full & ex-0DTE & p-value \\
\midrule
2020 & \tbd{} & \tbd{} & \tbd{} & \tbd{} & \tbd{} & \tbd{} \\
2021 & \tbd{} & \tbd{} & \tbd{} & \tbd{} & \tbd{} & \tbd{} \\
2022 & \tbd{} & \tbd{} & \tbd{} & \tbd{} & \tbd{} & \tbd{} \\
2023 & \tbd{} & \tbd{} & \tbd{} & \tbd{} & \tbd{} & \tbd{} \\
2024 & \tbd{} & \tbd{} & \tbd{} & \tbd{} & \tbd{} & \tbd{} \\
2025 & \tbd{} & \tbd{} & \tbd{} & \tbd{} & \tbd{} & \tbd{} \\
2026 & \tbd{} & \tbd{} & \tbd{} & \tbd{} & \tbd{} & \tbd{} \\
\bottomrule
\end{tabular}
\end{table}

\subsection{Deep Baseline, 2008--2019}
\label{sec:results:deep}

Figure~\ref{fig:deep} extends the analysis to 2008--2019, before the growth of weekly and daily expirations. \tbd{}
% TODO(results/deep): classifier and model rates by year for SPY on the three complete/diagnostic series; mark the 2008-2009 financial
% crisis and the 2018 volatility episode in the figure; report the gap between new_full and new_ex0dte over time as a consistency
% check of the ablation arm.

\begin{figure}[H]
\centering
\resultfigure{\figDeep}
\caption{Pre-0DTE baseline, SPY, 2008--2019, with the 2008--2009 financial crisis and the 2018 volatility episode marked. (a) Detection rate by year of the classifier (lines) and gpt-6-luna (markers). (b) Gap between \ident{new_full} and \ident{new_ex0dte}. (c) Share of $|\mathrm{GEX}|$ carried by same-day expiries. \textbf{Read this figure as:} \tbd{}}
\label{fig:deep}
\end{figure}

\subsection{Cross-Asset Replication and the Floor Effect}
\label{sec:results:cross}

Table~\ref{tab:cross} and Figure~\ref{fig:cross} repeat the analysis for QQQ and IWM. The fixed dollar thresholds of Section~\ref{sec:methods:criteria} can produce a floor effect for smaller underlyings, which we report; the scaled series of Section~\ref{sec:methods:arms} show whether the model's verdict still follows the criteria when the floor is lifted. \tbd{}
% TODO(results/cross): the data audit found zero persistent windows for QQQ and IWM in every year under the fixed $5B/$3B thresholds
% (floor effect); report it as a pre-registered design decision and use regime_detected as the primary label. The scaled-input variant
% (arms full_qqq_scaled and full_iwm_scaled on new_full_scaled, factors from series_metadata) was decided on 2026-10-28 and is reported
% next to the unscaled arms as a disclosed protocol change; a scaled-threshold variant would change the prompt text and stays out of
% scope. QQQ starts in 2011 (first window end 2011-05-04); the scaled series start in 2020.

\begin{table}[H]
\centering
\caption{Cross-asset replication: windows, typical magnitude and persistent windows under the fixed thresholds, by symbol and series. The scaled series are the disclosed protocol change of Section~\ref{sec:methods:arms}.}
\label{tab:cross}
\footnotesize
\begin{tabular}{@{}llrrrr@{}}
\toprule
Symbol & Series & Windows & \shortstack[r]{Median mean\\ $|\mathrm{GEX}|$ (\$B)} & \shortstack[r]{Persistent\\ (classifier)} & \shortstack[r]{Detected\\ (gpt-6-luna)} \\
\midrule
SPY & \ident{new_full} & \tbd{} & \tbd{} & \tbd{} & \tbd{} \\
QQQ & \ident{new_full} & \tbd{} & \tbd{} & \tbd{} & \tbd{} \\
IWM & \ident{new_full} & \tbd{} & \tbd{} & \tbd{} & \tbd{} \\
QQQ & \ident{new_ex0dte} & \tbd{} & \tbd{} & \tbd{} & \tbd{} \\
IWM & \ident{new_ex0dte} & \tbd{} & \tbd{} & \tbd{} & \tbd{} \\
QQQ & \ident{new_full_scaled} & \tbd{} & \tbd{} & \tbd{} & \tbd{} \\
IWM & \ident{new_full_scaled} & \tbd{} & \tbd{} & \tbd{} & \tbd{} \\
\bottomrule
\end{tabular}
\end{table}

\begin{figure}[H]
\centering
\resultfigure{\figCross}
\caption{Cross-asset replication, SPY, QQQ and IWM (QQQ from 2011). Top row: detection rate by year of the classifier (lines) and gpt-6-luna (markers) on \ident{new_full} and \ident{new_ex0dte}, and the share of windows that pass the persistence and sign-flip tests alone (dotted). Bottom row: yearly median and largest 30-day mean $|\mathrm{GEX}|$ against the \$5 billion criterion. \textbf{Read this figure as:} \tbd{}}
\label{fig:cross}
\end{figure}

\subsection{Post-Publication Period}
\label{sec:results:post}

Table~\ref{tab:post} and Figure~\ref{fig:post} cover the post-publication period of Section~\ref{sec:methods:windows}: the windows that end in 2026, of which those with ends from 2026-02-13 lie entirely after the end of the published sample (2025-12-31). We use this label rather than out-of-sample, because gpt-6-luna was created on 2026-09-14 and its training cutoff is undocumented. \tbd{}
% TODO(results/post): SPY windows ending in 2026 (to 2026-10-09) on new_full and new_ex0dte, with the counts for all of them and for those entirely in 2026 (ends from 2026-02-13). Explain the cap
% at 2026-10-09 (vendor switch on 2026-10-12) and that the model's training data may overlap the period.

\begin{table}[H]
\centering
\caption{Post-publication period, SPY, windows ending in 2026 up to 2026-10-09 (those with ends from 2026-02-13 lie entirely in 2026).}
\label{tab:post}
\footnotesize
\begin{tabular}{@{}lrrrr@{}}
\toprule
Series & Windows & Classifier detects (\%) & gpt-6-luna detects (\%) & Agreement (\%) \\
\midrule
\ident{new_full} & \tbd{} & \tbd{} & \tbd{} & \tbd{} \\
\ident{new_ex0dte} & \tbd{} & \tbd{} & \tbd{} & \tbd{} \\
\bottomrule
\end{tabular}
\end{table}

\begin{figure}[H]
\centering
\resultfigure{\figPost}
\caption{Post-publication period (windows ending in 2026 up to 2026-10-09): per-window classifier and gpt-6-luna labels for SPY under the complete series; the windows that still contain days of 2025 are set apart. \textbf{Read this figure as:} \tbd{}}
\label{fig:post}
\end{figure}

\subsection{Intraday Check}
\label{sec:results:intraday}

Table~\ref{tab:intraday} and Figure~\ref{fig:intraday} compare the late-session intraday snapshot with the end-of-day value for SPY from 2026-07-27 to 2026-10-09. \tbd{}
% TODO(results/intraday): snapshot type final_rush (19:55:59 UTC), the last independent late-session observation; the market_close slot
% equals the end-of-day row. Exclude 2026-07-29, 2026-08-03 and 2026-09-14 (irregular days); equity intraday rows before
% 2026-07-25 carried the previous day's chain (upstream #64). Compare composite windows (new_full for T-29..T-1 and the
% intraday value at T+0) with the end-of-day windows. With a few dozen days report counts, not p-values.

\begin{table}[H]
\centering
\caption{Late-session intraday snapshot versus end-of-day GEX, SPY, 2026-07-27 to 2026-10-09.}
\label{tab:intraday}
\footnotesize
\begin{tabular}{@{}lr@{}}
\toprule
Statistic & Value \\
\midrule
Days with both observations & \tbd{} \\
Days with the same sign (\%) & \tbd{} \\
Median relative difference of GEX (\%) & \tbd{} \\
Classifier label agreement of composite and end-of-day windows (\%) & \tbd{} \\
gpt-6-luna label agreement of composite and end-of-day windows (\%) & \tbd{} \\
\bottomrule
\end{tabular}
\end{table}

\begin{figure}[H]
\centering
\resultfigure{\figIntraday}
\caption{Late-session intraday snapshot versus end-of-day GEX, SPY, 2026-07-27 to 2026-10-09. (a) Day T+0 value of the snapshot (\ident{final_rush}) and of the end-of-day row. (b) 30-day mean $|\mathrm{GEX}|$ of the composite and of the end-of-day window against the \$5 billion criterion. (c) Labels of the classifier and of gpt-6-luna on the two windows. Irregular days are drawn hollow and not counted. \textbf{Read this figure as:} \tbd{}}
\label{fig:intraday}
\end{figure}

\subsection{Model Fidelity to the Rule}
\label{sec:results:fidelity}

Table~\ref{tab:fidelity} and Figure~\ref{fig:fidelity} compare each model answer with the deterministic classifier on the same window. \tbd{}
% TODO(results/fidelity): accuracy of regime_detected against clf_should_detect by arm, Cohen's kappa, the share of answers whose
% reported counts, magnitude and persistence equal the computed values (metrics_agree: 0.5 day, 0.05B, 0.5 percentage points),
% the exact-match rate of sign_flips (undercounted in 7 of 32 in the probe), and behaviour in the borderline region around 70 %
% persistence and 5B magnitude.

\begin{table}[H]
\centering
\caption{Fidelity of model answers to the deterministic classifier, by arm (the published o4-mini answers of \ident{r1_o4mini} are the reference).}
\label{tab:fidelity}
\footnotesize
\begin{tabular}{@{}lrrrrr@{}}
\toprule
Arm & Windows & Accuracy (\%) & Cohen's kappa & Metrics agree (\%) & Flips exact (\%) \\
\midrule
\ident{r1_o4mini} & \tbd{} & \tbd{} & \tbd{} & \tbd{} & \tbd{} \\
\ident{bridge_jan} & \tbd{} & \tbd{} & \tbd{} & \tbd{} & \tbd{} \\
\ident{full_spy} & \tbd{} & \tbd{} & \tbd{} & \tbd{} & \tbd{} \\
\ident{ex0dte_spy} & \tbd{} & \tbd{} & \tbd{} & \tbd{} & \tbd{} \\
\ident{full_qqq} & \tbd{} & \tbd{} & \tbd{} & \tbd{} & \tbd{} \\
\ident{full_iwm} & \tbd{} & \tbd{} & \tbd{} & \tbd{} & \tbd{} \\
\ident{full_qqq_scaled} & \tbd{} & \tbd{} & \tbd{} & \tbd{} & \tbd{} \\
\ident{full_iwm_scaled} & \tbd{} & \tbd{} & \tbd{} & \tbd{} & \tbd{} \\
\ident{p34_obf_2024} & \tbd{} & \tbd{} & \tbd{} & \tbd{} & \tbd{} \\
\ident{p34_dated_2024} & \tbd{} & \tbd{} & \tbd{} & \tbd{} & \tbd{} \\
\ident{p34_obf_2020} & \tbd{} & \tbd{} & \tbd{} & \tbd{} & \tbd{} \\
\ident{p34_dated_2020} & \tbd{} & \tbd{} & \tbd{} & \tbd{} & \tbd{} \\
\bottomrule
\end{tabular}
\end{table}

\begin{figure}[H]
\centering
\resultfigure{\figFidelity}
\caption{Fidelity of the model answers (o4-mini and gpt-6-luna) to the deterministic classifier. (a), (b) Accuracy and Cohen's kappa of the verdict by series and year. (c)--(e) Reported versus computed persistence, mean $|\mathrm{GEX}|$ and sign-flip count, with the borderline region around each threshold shaded. \textbf{Read this figure as:} \tbd{}}
\label{fig:fidelity}
\end{figure}

\subsection{Dated versus Obfuscated Prompts: The 2024 and 2020 Baselines}
\label{sec:results:dated}

Table~\ref{tab:dated} compares the published o4-mini answers on the inputs of the Phase 3 and 4 baselines, whose prompts carried calendar dates, with gpt-6-luna on the same values with the dates and with relative day labels (Section~\ref{sec:methods:arms}). \tbd{}
% TODO(results/dated): arms r1_p3_o4mini and r1_p4_o4mini (published o4-mini answers, dated prompts; data audit: 181 of 223 = 81.2 % for 2024
% and 27 of 223 = 12.1 % for 2020, to be verified against the final tables), p34_dated_* and p34_obf_*: detection rate with Wilson
% intervals, agreement with the classifier on the series as sent (p3_2024 and p4_2020), the paired dated-versus-obfuscated difference per
% baseline (exact McNemar on the non-overlapping subsample, quarter-block bootstrap) and the 2024-versus-2020 contrast (phi; 0.692 for the
% published answers) under each label type. The two series differ in sign convention and in size (mean |GEX| about 31.2 B against 2.85 B in
% the data audit), so say that the contrast is one of magnitudes and do not describe it as negative versus positive regimes. The 223
% windows of each baseline overlap (stride one): report the non-overlapping subsample next to every p-value.

\begin{table}[H]
\centering
\caption{Phase 3 (2024) and Phase 4 (2020) baselines: detection rate in percent (detected windows out of 223) of the published o4-mini answers, which carried dates, of gpt-6-luna with the dates and with relative day labels (Day T), and of the deterministic classifier on the series as sent. The last column is the exact McNemar p-value of the two gpt-6-luna arms on the non-overlapping subsample (Section~\ref{sec:methods:statistics}).}
\label{tab:dated}
\footnotesize
\begin{tabular}{@{}llrrrrr@{}}
\toprule
Baseline & Series & o4-mini, dated & luna, dated & luna, Day T & Classifier & McNemar p \\
\midrule
Phase 3, 2024 & \ident{p3_2024} & \tbd{} & \tbd{} & \tbd{} & \tbd{} & \tbd{} \\
Phase 4, 2020 & \ident{p4_2020} & \tbd{} & \tbd{} & \tbd{} & \tbd{} & \tbd{} \\
\bottomrule
\end{tabular}
\end{table}
